\begin{lemma}
\label{lemma-silly-normal}
Let $R$ be a ring. Let $x, y \in R$ be nonzerodivisors.
Let $R[x/y] \subset R_{xy}$ be the $R$-subalgebra generated
by $x/y$, and similarly for the subalgebras $R[y/x]$ and $R[x/y, y/x]$.
If $R$ is integrally closed in $R_x$ or $R_y$, then the sequence
$$
0 \to R \xrightarrow{(-1, 1)} R[x/y] \oplus R[y/x] \xrightarrow{(1, 1)}
R[x/y, y/x] \to 0
$$
is a short exact sequence of $R$-modules.
\end{lemma}

\begin{proof}
Because $x$ and $y$ are nonzerodivisors, their product is a
nonzerodivisor. Thus $R\to R_{xy}$ is injective, and so are the
maps $R_x\to R_{xy}$ and $R_y\to R_{xy}$: a fraction that becomes
zero has its numerator killed by a power of $xy$, hence has zero
numerator. All subrings in the statement therefore have their
specified inclusions in $R_{xy}$. The product of the original
elements $x/y$ and $y/x$ is $1$. Consequently each monomial
$(x/y)^a(y/x)^b$ with $a,b\geq0$ equals $(x/y)^{a-b}$ when
$a\geq b$, and equals $(y/x)^{b-a}$ otherwise. Splitting a
polynomial into these terms expresses it as a sum of elements
of $R[x/y]$ and $R[y/x]$, proving surjectivity of the map
$(u,v)\mapsto u+v$. Its kernel consists exactly of pairs
$(-\alpha,\alpha)$ with $\alpha\in R[x/y]\cap R[y/x]$.
The first map in the statement is injective and sends
$a\mapsto(-a,a)$, with the original signs.
Let $\alpha \in R[x/y] \cap R[y/x]$. To show exactness in the middle
we have to prove that $\alpha \in R$. By assumption we may write
$$
\alpha = a_0 + a_1 x/y + \ldots + a_n (x/y)^n =
b_0 + b_1 y/x + \ldots + b_m(y/x)^m
$$
for some $n, m \geq 0$ and $a_i, b_j \in R$.
Pick some $N > \max(n, m)$.
Consider the finite $R$-submodule $M$ of $R_{xy}$ generated by the elements
$$
(x/y)^N, (x/y)^{N - 1}, \ldots, x/y, 1, y/x, \ldots, (y/x)^{N - 1}, (y/x)^N
$$
We check $\alpha M\subset M$ using the full original expressions.
For $0\leq i\leq N$, multiply $(x/y)^i$ by
$\sum_{j=0}^m b_j(y/x)^j$. If $i\geq j$, the resulting
monomial is $(x/y)^{i-j}$ with $0\leq i-j\leq N$; if $i<j$,
it is $(y/x)^{j-i}$ with $0<j-i\leq m<N$.
Thus every summand, with its original coefficient $b_j$, belongs
to $M$. Multiplying $(y/x)^i$ instead by
$\sum_{j=0}^n a_j(x/y)^j$ gives $(y/x)^{i-j}$ for $i\geq j$
and $(x/y)^{j-i}$ for $i<j$, with the same bounds using $n<N$.
This proves stability for every listed generator. Since $1\in M$,
Lemma \ref{lemma-characterize-integral-element} proves $\alpha$
integral over $R$. The expression with coefficients $b_j$ places
$\alpha$ in the original $R_x$, and the expression with $a_j$
places it in $R_y$. Either of the two integral-closedness
hypotheses therefore gives $\alpha\in R$, proving middle exactness.

The calculation also identifies precisely what is needed without
either integral-closedness assumption. With the same original
nonzerodivisors, let $C=R[x/y]\cap R[y/x]\subset R_{xy}$.
It is a subring containing $R$, and the preceding finite-submodule
argument proves every element of $C$ integral over $R$.
There is always an exact sequence with its actual intersection
$$
0\longrightarrow C\xrightarrow{\alpha\mapsto(-\alpha,\alpha)}
R[x/y]\oplus R[y/x]\xrightarrow{(u,v)\mapsto u+v}
R[x/y,y/x]\longrightarrow0.
$$
Surjectivity and the kernel were proved above, and injection uses
the original inclusions. In the sequence with $R$ at the left,
the middle homology is exactly $C/R$: the map sends
$\alpha+R$ to the class of $(-\alpha,\alpha)$, is well defined
by the original map $R\to R[x/y]\oplus R[y/x]$, and has inverse
given by the second coordinate of a kernel pair modulo $R$.
Both composites are identities. Thus the original sequence is
exact precisely when this specified integral extension $C$ equals
$R$; the two stated hypotheses are sufficient conditions for that
equality, and the formula retains the precise defect otherwise.
\end{proof}